\section{Implementation and Development}
The development of the Smart Plant Watering System followed an iterative approach, transitioning from functional simulation to breadboard prototyping and finally to the stylized diorama implementation.

\subsection{Development Environment}
The project was developed using the \textbf{MPLAB X IDE} with the \textbf{XC8 v2.45} compiler. Programming and debugging were performed using the \textbf{MPLAB Snap} debugger.

\subsection{Challenges and Solutions}
\subsubsection{ADC Noise and Stability}
Initial tests showed significant fluctuations in moisture readings when the pump was active. This was traced back to the electrical noise generated by the DC motor affecting the common ground.
\begin{itemize}
    \item \textbf{Solution:} Implementation of the 100-sample averaging filter and the addition of a decoupling capacitor near the sensor's power pins. Furthermore, the software-controlled sensor power (\texttt{HYGRO\_POWER\_PIN}) significantly reduced baseline noise.
\end{itemize}

\subsubsection{I2C Timing Constraints}
Implementing I2C via bit-banging required precise delay management to respect the standard-mode (100kHz) timing.
\begin{itemize}
    \item \textbf{Solution:} Fine-tuning the \texttt{\_\_delay\_us()} calls within the I2C driver to ensure reliable communication with the LCD controller (PCF8574) without blocking the main execution loop for too long.
\end{itemize}

\subsection{Miniaturization and Aesthetics}
One of the unique aspects of this project was the creation of a "Smart Garden" diorama. The challenge was to hide the electronics (breadboard, wires, and microcontroller) within a miniature wooden cabin while keeping the LCD visible and accessible. This required careful wire routing and the use of extensions for the sensor and pump.
